\scanpage{48}
\mspage{\handwritten{19}}
% The typescript is upside down in the scan. A long pen stroke cancels the sheet.
\deleted{9.1.13. \uline{Questions ouvertes}. Soit $E$ un \uline{U}-topos.}

\deleted{a) Dans l'ensemble \inserted{ordonné} des sous-topos de\unsure{,} $E$, le Inf de deux éléments est-il distributif par rapport aux Sup finis ? Pour des résultats dans c e \deleted{\unsure{se|sens}} sens, voir 9.1.11 e) et 9.1.12 d) .}

\deleted{b) Soit $E'$ un sous-topos de $E$, distinct de \deleted{sous-topos ``vide''} $E$. Existe-t-il un sous-topos $E''$ de $E$, tel que $E''\handwritten{\cap}E'=\operatorname{Top}(\varnothing)$, $E''\handwritten{\ne}\operatorname{Top}(\varnothing)$ ?}

\deleted{\deleted{\unsure{b)|p}}c) Soient $E'$ et $E''$ deux sous-topos de $E$ tels que $E'\handwritten{\cap}E''=\operatorname{Top}(\varnothing)$. Les conditions suivantes sont-elles équivalentes: (i) $E$ est le sup de $E'$ et $E''$ dans l'ensemble des sous-topos de $E$ (on se retient à grand peine d'écrire cette condition $E=E'\handwritten{\deleted{\unsure{\cap|\vee}}\cup}E''$ ); (ii) $E''$ est \deleted{max} la plus grand parmi les sous-topos $E''_{1}$ de $E$ tels que $\operatorname{Top}(E')\handwritten{\cap}\operatorname{Top}(E''_{1})=\operatorname{Top}(\varnothing)$; (iii) $E''$ est maximal parmi les $E''_{1}$ précédents.}

\deleted{Il est \deleted{\unsure{a bit|abit|abit.}} clair que (ii)$\handwritten{\Longrightarrow}$(iii), que si a) est vérifié \unsure{en|on} a (i)$\handwritten{\Longrightarrow}$(ii), \deleted{\unsure{et que si|e que si}} \inserted{que} b) est \unsure{vérifié|vérifi} \deleted{\unsure{e|é}} \inserted{si et seulement si \deleted{on a}} on a (iii)$\handwritten{\Longrightarrow}$(i); donc une solution affirmative à a) et b) en impliquerait une pour c). Noter aussi que la réponse aux questions posées \inserted{( y compris d) qui suit)} est affirmative si $E$ est de la forme $\operatorname{Top}(X)$, $X$ un espace topologique \inserted{localement} noethérien (9.1.11 b)), \deleted{ainsi que pour \unsure{$E=G$|$E\in G$}} \deleted{ou si $E$ est} un topos fini (9.1.12).}

\deleted{d) Si $E$ est le sup des deux sous-topos $E'$ et $E''$, alors tout point de $E$ est-il isomorphe à un point provenant de $E'$ ou de $E''$ ?}

\deleted{Si la réponse à d) était affirmative, alors on aurait une réponse affirmative à a) \inserted{$(E'\ (E''\ E''')=E'\ E''\quad E'\ E''')$} pour le\deleted{s} cas de sous-topos $E',E'',E'''$ de $E$ ayant suffisamment de points (donc a) serait vrai pour $E$ \inserted{localement} noethérien, en vertu de 9.1.11 b)). Inversement, si a)}
% archive sheet number: 47
